\chapter{Postprocessors}\label{ch:posts}

\begin{warningbox}
Every postprocessor can crash a machine, the ones that come with WOP too. None of them was run on every controller. Read the G-code, and cut in air first, before you run a new postprocessor on a part.
\end{warningbox}

\section{What a postprocessor does}
WOP plans the paths in its own form: moves, spindle, coolant, tool changes. A \emph{postprocessor} (post) turns them into the G-code of your controller. It decides how a rapid move is written, whether X is a diameter, how a thread or a drill cycle looks, what the first lines say, and so on.

A post is a text file that ends in \file{.post.ini}. You do not need a programmer to change it. You change it in the postprocessor editor, which shows what the G-code will look like while you edit.

One post comes with WOP: the post for \textbf{LinuxCNC}. It writes X as diameter, feed per revolution, constant surface speed, the thread cycle and the drill cycles of LinuxCNC, and cutter compensation (G41, G42, G40) for a finishing pass that asks for it. It was tried on one lathe.

Inside the program there is also a plain default post (plain G0, G1, G2, G3 and a few M codes, feed in mm/min, no cycles). It is not in the lists. A post that does not have a page of its own for something uses the page of this default post.

The folder \file{examples/posts} holds reference posts to import with \button{Import...}: a Fanuc style post, a plain G-code post, and the LinuxCNC post with a custom machine function (a bar feeder) that shows how a custom function is written. The Fanuc and the plain post are written from a rough understanding of these controllers and are \emph{not tried} on any machine. Use them as a start for your own post.

The machine says which post is used (Chapter~\ref{ch:machines}).

\section{The editor}\label{sec:posts:editor}
Open it with \menupath{File}{Postprocessors...} (Figure~\ref{fig:post_editor}).

\shot[0.9]{post_editor}{The postprocessor editor.}

On top is the list of the posts and the buttons:
\begin{description}
\item[\button{New}] Makes a new post.
\item[\button{Save}] Saves your changes.
\item[\button{Save as copy...}] Saves under a new name. The posts that come with WOP cannot be changed; make a copy and change that.
\item[\button{Import...}] Copies a post file of somebody else into your posts.
\item[\button{Delete}] Removes one of your posts.
\item[\button{Folder}] Shows the folder of your posts in Windows.
\item[\field{Sample X}] Whether the sample on the pages shows X as diameter or radius.
\end{description}
The title of the window gets a star when there are unsaved changes.

\section{The pages of a post}\label{sec:posts:pages}
The pages are on the left. They are in this order.

\begin{description}
\item[Post] The name and the description of the post, and the name of its file. A note says whether the whole file reads without a problem.
\item[Settings] Numbers and switches: the number of decimals, line numbers, whether X is a diameter, whether arcs can be written while X is a diameter, how a thread and a drill are written (as cycle, in two blocks, or as plain passes), cutter compensation, a speed ramp before the spindle stops (for a drive that trips when it is stopped from speed), the characters a comment may not hold, and the file extension. Hover over a name to see what it means.
\item[Code table] The G and M codes of the controller, by name. WOP uses the names; you set the codes. A code that is empty means the controller has no such function, so WOP does not write it. The section \emph{Work offsets} is the list of work offsets of the controller (G54, G55 \ldots) that the Setup offers; add the ones your controller has. You can add codes of your own.
\item[Modal groups] Which codes cancel each other (for example G0, G1, G2 and G3 are one group). WOP uses the groups to leave out a code that is still valid from the line before. List the codes of the code table here; a code of the table that is in no group is named on the Code table page. You can add groups of your own.
\item[The concepts] After these four pages comes one page for each thing the post writes: the header and the footer, comments, the start of a block, rapid, the move to the tool change position (a page of its own, for a controller that has a code for it, such as G28 U0 W0), feed, arcs, the thread cycle, the drill cycles, a dwell, stops, the tool change, the spindle, coolant, and more. A page has one \emph{template}. Pages that the post does not have of its own show a note: the post uses the default post for this. Tick the box to make it the post's own; it starts as a copy.
\end{description}

\subsection{Writing a template}
A template is the line or lines that the post writes for this thing. It is a text with \emph{tokens} between angle brackets. A token is replaced by a value: a code, a coordinate, a feed. For example:
\begin{center}\texttt{[<comp code> ]<code feed changed>[ X<X changed>][ Z<Z changed>][ F<F changed>]}\end{center}
is the template of a feed move in the LinuxCNC post. It writes \texttt{G1 X20 Z-5 F0.15}, and leaves out each value that did not change. What the pieces mean:
\begin{itemize}
  \item \texttt{<code feed>} is the code with the name \emph{feed} from the code table; \texttt{<X>} is the X coordinate. A token can have modifiers; \texttt{changed} writes it only when it differs from the line before.
  \item Square brackets make a part \emph{optional}. If the token inside has nothing to write, the whole part (including the space before the X) is left out.
\end{itemize}
Below the template, the page shows the same template written twice: with every optional part present and everything changed, and with nothing changed and every optional part gone. Look at both. A problem in the template is written in red.

Under the template is the list of tokens that work on this page. Double click a token, or click it, to put it into the template. \button{All tokens...} lists all of them.

\section{Making your own post}\label{sec:posts:own}
\begin{enumerate}
  \item In the editor, choose the post that is nearest to your controller and press \button{Save as copy...}. Give it a name.
  \item Go through the \emph{Settings} and the \emph{Code table}. Most differences between controllers are here.
  \item Go through the concept pages that matter: the tool change, the spindle and the cycles are the usual ones. Tick \emph{own} on a page only when you need to change it.
  \item Press \button{Save}. In \menupath{File}{Machines...}, choose the post in the machine.
\end{enumerate}
A post of your own is kept in the \file{posts} folder of the user folder (Chapter~\ref{ch:app_files}). You can copy the file to another PC.

\section{Checking what a post writes}\label{sec:posts:check}
Do this every time you make or change a post, before the machine runs it:
\begin{enumerate}
  \item Export the G-code of a small project that uses the things you changed (Section~\ref{sec:file:export}).
  \item Open the file in a text editor. Read it from the top. Check the first lines (units, plane, absolute or incremental), the tool change, the spindle start, one rapid move and one feed move.
  \item Run it on the machine without the part and without the tool, or at a safe height, with your hand on the stop.
\end{enumerate}
